% !TEX root = ../main.tex
\chapter*{KATA PENGANTAR}
\addcontentsline{toc}{chapter}{KATA PENGANTAR}

Puji dan syukur penulis panjatkan ke hadirat Tuhan Yang Maha Esa atas rahmat dan karunia-Nya sehingga laporan Tugas Akhir berjudul ``Pemanfaatan Generative AI untuk Process Discovery pada Data Operasional Pemeliharaan Perusahaan Listrik'' dapat diselesaikan dengan baik. Laporan ini disusun sebagai salah satu syarat memperoleh gelar sarjana pada Program Studi Sarjana Informatika, Fakultas Informatika, Universitas Telkom.

Tugas Akhir ini membahas pemanfaatan model bahasa untuk menstandarkan ragam label aktivitas pemeliharaan sebelum model proses dibangun, sekaligus membandingkan hasilnya dengan pendekatan \textit{process discovery} konvensional pada data operasional pemeliharaan pembangkit listrik.

Penulis menyadari laporan ini masih memiliki kekurangan. Karena itu, penulis mengharapkan kritik dan saran yang membangun demi perbaikan di masa mendatang, dengan harapan laporan ini dapat bermanfaat bagi pembaca.

\vspace{1cm}

\begin{flushright}
  Bandung, 07 April 2026 \\[2cm]
  \textbf{Sakha Wibisono}
\end{flushright}

\chapter*{UCAPAN TERIMA KASIH}
\addcontentsline{toc}{chapter}{UCAPAN TERIMA KASIH}

Penulis menyampaikan terima kasih yang sebesar-besarnya kepada seluruh pihak yang telah memberikan bantuan, bimbingan, dan dukungan selama penyusunan Tugas Akhir ini, khususnya kepada:

\begin{enumerate}
  \item Ibu Angelina Prima Kurniati, S.T., M.T., Ph.D., selaku dosen pembimbing, atas arahan, bimbingan, dan masukan yang diberikan sepanjang proses penelitian.
  \item Bapak Mahmud Dwi Sulistiyo, S.T., M.T., Ph.D., selaku Ketua Program Studi Sarjana Informatika, beserta seluruh dosen dan staf Fakultas Informatika Universitas Telkom, atas ilmu dan bantuan yang diberikan selama masa perkuliahan.
  \item PT PLN (Persero) selaku pengelola unit PLTD Titikuning, atas kesediaannya menyediakan data operasional pemeliharaan yang menjadi dasar penelitian ini.
  \item Kedua orang tua penulis, Bapak Suriyanto dan Ibu Tuti Afitri, serta kakak penulis, Shata Pramudya Ananda, dan adik penulis, Maisah Azzura, beserta segenap keluarga, atas doa, dukungan, dan semangat yang tidak pernah putus.
  \item Rekan mahasiswa, sahabat, serta semua pihak yang tidak dapat penulis sebutkan satu per satu, atas kebersamaan dan dukungannya.
\end{enumerate}

Semoga segala kebaikan dan bantuan yang telah diberikan mendapat balasan yang setimpal dari Tuhan Yang Maha Esa.

